\begin{annexitem}{annex:p19.r33.witness}{thm:p19-r33-witness}{Witnesses and the difficult source inference}
\begin{proof}
For \(Iab,Abc\), choose \(z\in a^U\cap b^U\).
Since \(b^U\subseteq c^U\), the same \(z\) belongs to
\(a^U\cap c^U\). This proves \(Iac\).
For \(Eab,Abc\) with nonempty \(b^U\), choose \(z\in b^U\).
It belongs to \(c^U\) by inclusion and not to \(a^U\) by
disjointness. Hence \(z\in c^U\setminus a^U\), proving \(Oca\).
To refute \(Eac\) under the same conditions, take
\(U=\{d,m\}\), \(a^U=\{d\}\), \(b^U=\{m\}\),
\(c^U=\{d,m\}\). The premises and all three nonemptiness
conditions hold, whereas \(a^U\cap c^U\) contains \(d\).
Finally take \(U=\{d\}\), \(a^U=c^U=\{d\}\),
\(b^U=\varnothing\). Both premises remain true but
\(c^U\setminus a^U=\varnothing\). Thus neither the universe's
nonemptiness nor the end classes' nonemptiness repairs the missing
middle-class witness.
\end{proof}
\end{annexitem}

\begin{annexitem}{annex:p19.r33.nvc}{thm:p19-r33-nvc}{Countermodels for every requested end-term answer}
\begin{proof}
For \(Aab,Obc\), use the following three complete interpretations.
Every unlisted property membership is false.
\[
\begin{array}{c|c|c|c|c}
&U&a^U&b^U&c^U\\
1&\{d,w,e\}&\{d\}&\{d,w\}&\{d,e\}\\
2&\{d,w,e\}&\{d\}&\{d,w\}&\{e\}\\
3&\{d,w\}&\{d\}&\{d,w\}&\{d\}.
\end{array}
\]
All three have nonempty named classes, \(a^U\subseteq b^U\)
and a witness \(w\in b^U\setminus c^U\).
Interpretation 1 refutes \(Eac,Eca,Oac,Aca\): it has end-class
overlap, no \(a\)-member outside \(c\), and \(e\in c^U\setminus a^U\).
Interpretation 2 refutes \(Aac,Iac,Ica\), since \(d\notin c^U\)
and the end classes are disjoint. Interpretation 3 refutes \(Oca\).
Together they refute each of the eight proposed answers while retaining
the premises, establishing NVC for the specified answer space.

For \(Aba,Ocb\), use
\[
\begin{array}{c|c|c|c|c}
&U&a^U&b^U&c^U\\
1&\{x,y\}&\{x\}&\{x\}&\{y\}\\
2&\{x,y\}&\{x,y\}&\{x\}&\{y\}\\
3&\{x,y\}&\{x\}&\{x\}&\{x,y\}.
\end{array}
\]
Each has \(b^U\subseteq a^U\), a witness
\(y\in c^U\setminus b^U\), and nonempty classes.
Interpretation 1 refutes \(Aac,Aca,Iac,Ica\).
Interpretation 2 refutes \(Eac,Eca,Oca\), and interpretation 3
refutes \(Oac\). This covers every requested sentence.
The premises themselves remain logical consequences, so the result
does not assert absence of all consequence.
\end{proof}
\end{annexitem}

\begin{annexitem}{annex:p19.r33.split}{thm:p19-r33-split}{One merged witness versus two distinct witnesses}
\begin{proof}
On \(U=\{z\}\), let \(a^U=b^U=c^U=\{z\}\).
Then both existential premises and \(Iac\) hold. Split the triple
membership into two individuals on \(U'=\{x,y\}\), with
\(a^{U'}=\{x\}\), \(b^{U'}=\{x,y\}\),
\(c^{U'}=\{y\}\). The first premise has witness \(x\),
the second witness \(y\), and \(a^{U'}\cap c^{U'}=\varnothing\).
Thus the two premises also have a model satisfying \(Eac\).
Neither answer follows in all models. Every named extension is nonempty
in both interpretations, so existence assumptions do not rescue the
matching inference. The construction demonstrates a possible revision,
not a claim that any participant actually performed it.
\end{proof}
\end{annexitem}

\begin{annexitem}{annex:p19.r33.sampling}{thm:p19-r33-sampling}{The finite-sample boundary and test exhaustiveness}
\begin{proof}
Each draw misses the specified overlap with probability \(1-q\).
Independence gives the product \((1-q)^n\); for \(n=0\), the
empty product is 1. If \(0<q<1\) and \(n\) is finite, every
factor is positive, so the product is positive. The population or
prototype can have overlap while a finite sample contains none.
For a fixed sample, \(A\) means no sampled member of \(S\)
lies outside \(P\), while \(O\) means one does. Exactly one
is true. Similarly, \(I\) means a sampled overlap exists and
\(E\) that none exists; exactly one is true. This includes an
empty \(S\) extension: \(A,E\) then pass vacuously.
Consequently an algorithm inspecting all four and accepting the first
passing test must find one. With a gate rejecting a passing candidate
with positive probability, that reasoning no longer establishes an
accepted output. Such gates, present in the full Hattori algorithm,
can yield NVC without a counterexample to every premise-level candidate.
\end{proof}
\end{annexitem}

\begin{annexitem}{annex:p19.r33.finite-check}{thm:p19-r33-finite-check}{Exhausting membership types}
\begin{proof}
An individual has one of the eight complete membership triples for
\(a,b,c\). Given any nonempty model, choose one representative of
each realized triple and form the substructure on those representatives.
There are at most eight and at least one. Every universal inclusion or
disjointness sentence true in the model remains true in the substructure.
Conversely, any violation of inclusion or disjointness has a witness of
a realized triple, whose representative preserves the violation.
Every existential intersection or difference also has a witness of a
realized triple, whose representative preserves it; absence of such a
witness is preserved as well. The same argument preserves nonemptiness
of each named class. Thus truth and falsity of every A/I/E/O sentence
are preserved, including all premises and the proposed conclusion.
Any countermodel therefore reduces to one of the nonempty subsets of
the eight types. Enumerate these 255 subsets, interpret membership by
their coordinates, discard those failing premises or declared existence
conditions, and inspect the proposed conclusion in the rest. If one
fails it, consequence is false; if none does, consequence holds by the
reduction. If no premise model exists, ordinary consequence is vacuous;
this is distinct from a satisfiable task with NVC. Repeating for the
eight candidates establishes the requested answer set. This finite
exhaustion is a formal benchmark, not an empirically established model
of human processing.
\end{proof}
\end{annexitem}
